\section*{第 \ref{ch:g3:food-chains} 章——食物链}

\begin{solution}{exo:g3:food-chains:1}
\omterm{def:g3:food-chains:chain}{“被……吃”}。箭头从被吃的一方指向吃它的一方——
顺着食物的方向。
\end{solution}

\begin{solution}{exo:g3:food-chains:2}
\omterm{def:g1:plants-around-us:plant}{植物}（或其他能自己制造食物的\omterm{def:g1:living-or-not:living}{生物}）。
因为\omterm{def:g1:plants-around-us:plant}{植物}只靠光、水和空气就能取食，不需要吃别的\omterm{def:g1:living-or-not:living}{生物}——
所有\omterm{def:g1:animals-around-us:animal}{动物}环节都依赖第一个绿色环节，
或直接、或间接地依靠它生活。
\end{solution}

\begin{solution}{exo:g3:food-chains:3}
草 $\rightarrow$ 蚱蜢 $\rightarrow$ 蜥蜴
$\rightarrow$ 鹰。
\end{solution}

\begin{solution}{exo:g3:food-chains:4}
箭头方向反了：草 $\rightarrow$ 兔子
$\rightarrow$ 狐狸——草被兔子吃，兔子
被狐狸吃。
\end{solution}

\begin{solution}{exo:g3:food-chains:5}
橡树叶 $\rightarrow$ \omterm{ex:g2:animals-grow-up:butterfly}{毛毛虫} $\rightarrow$ 蓝山雀
$\rightarrow$ 雀鹰。
\end{solution}

\begin{solution}{exo:g3:food-chains:6}
沙拉被蜗牛吃掉；刺猬吃蜗牛。
\end{solution}

\begin{solution}{exo:g3:food-chains:7}
例如：水生\omterm{def:g1:plants-around-us:plant}{植物} $\rightarrow$ \omterm{ex:g2:animals-grow-up:frog}{蝌蚪} $\rightarrow$
刺鱼 $\rightarrow$ 苍鹭。
\end{solution}

\begin{solution}{exo:g3:food-chains:8}
因为它的食物的食物的食物是草：褐色的草地上蚱蜢变少，
蜥蜴就少，
蜥蜴少了，鹰的捕食就变得困难。第一个环节的变化
会沿着整条\omterm{def:g3:food-chains:chain}{食物链}向上传递。
\end{solution}

\begin{solution}{exo:g3:food-chains:9}
毒死蜗牛，会让它们的捕食者挨饿：刺猬会离开
或死去。之后，蜗牛的主要天敌刺猬——
消失了，剩下的鼻涕虫和蜗牛繁殖得比以前
更快：去掉一个环节，影响既会沿\omterm{def:g3:food-chains:chain}{食物链}向上传递（饥饿的
刺猬），也会向下传递（没人捕食的蜗牛）。
\end{solution}

\begin{solution}{exo:g3:food-chains:10}
对——通过\omterm{def:g3:food-chains:chain}{食物链}中间的环节。蜥蜴吃蚱蜢，
蚱蜢吃草：如果草没有了，这种损失会
一环一环地向上传递，从草传到蚱蜢，再传到蜥蜴。
蜥蜴从不直接接触草，
却同样依赖草。
\end{solution}
